\section{A rank-uniform local shell construction}
\label{sec:wick}

This section is a separate result in its own regime, included because it
is exact, rank-uniform, and independent of everything above --- not
because it continues it. The object is the formal local Wilson
oscillator expansion around a single plaquette, and the point is that
its spectrum admits a finite exact recursion at every perturbative
order with coefficients in $\Q[N, N^{-1}]$, needing no Gram-matrix
inversion.

Conflating this construction with the flux band of
\S\ref{sec:carrier}--\S\ref{sec:sixth}, or with the Wilson transfer
operator of \S\ref{sec:wilson}, would be an error. They are four
different objects --- the physical $\SU(N)$ transfer band, the $\SU(3)$
internal Hamiltonian carrier, the conditional block, and this local
oscillator --- and no cross-regime continuation between them is
asserted anywhere in this document.

\subsection{The exact algebra}

Let $X$ be a traceless Hermitian $N \times N$ matrix with density
proportional to $e^{-\Tr(X^2)}$, write $P_k = \Tr(X^k)$ with $P_0 = N$
and $P_1 = 0$, and use a traceless Hermitian basis normalized by
$\Tr(T_a T_b) = \delta_{ab}$. The covariance is
\begin{equation}
  \mathbb{E}\big[X_{ab} X_{cd}\big]
  = \tfrac12\left(\delta_{ad}\delta_{bc}
    - \tfrac{1}{N}\delta_{ab}\delta_{cd}\right).
\end{equation}
With $D = \tfrac12\Delta$ and $E$ the polynomial Euler operator, the
completeness identity for this basis gives the exact relations
\begin{align}
  D P_k &= \frac{k}{2}\sum_{l=0}^{k-2} P_l P_{k-2-l}
           - \frac{k(k-1)}{2N} P_{k-2},
  \\
  D(fg) &= (Df)g + f(Dg) + \nabla f \cdot \nabla g,
  \\
  \nabla P_k \cdot \nabla P_l
        &= k\,l\left(P_{k+l-2} - \frac{1}{N} P_{k-1}P_{l-1}\right).
\end{align}
These are exact at every integer rank, including ranks at which some
traces become dependent --- which is the property that makes the
construction rank-uniform rather than rank-by-rank. For homogeneous $F$
of positive degree $d$, Gaussian integration by parts gives
$\mathbb{E}_\mu[F] = \mathbb{E}_\mu[DF]/d$, and with
$\mathbb{E}_\mu[1] = 1$ and vanishing odd moments this is a finite exact
moment recursion.

\subsection{Shell projectors without a Gram inverse}

\begin{proposition}[Terminating Wick conjugation]
\label{prop:wick}
Put $L = E - D$. Since $[E, D] = -2D$, the terminating exponential
$W = \exp(-D/2)$ satisfies
\begin{equation}
  L\,W = W\,E .
\end{equation}
Hence, with $H_d$ the extraction of the homogeneous degree-$d$ part,
\begin{equation}
  \Pi_d = \exp(-D/2)\, H_d \,\exp(D/2)
\end{equation}
is the exact oscillator-shell projector, and on each finite polynomial
space $R_s = \sum_{d \neq s} \Pi_d/(s-d)$ is the exact shell resolvent.
\end{proposition}

The content of Proposition~\ref{prop:wick} is that $L$ and $E$ are
conjugate by a \emph{terminating} exponential, so the shell
decomposition of the interacting operator is the shell decomposition of
the Euler operator, transported. No Gram matrix is formed and none is
inverted, which is precisely the step that makes rank-uniformity
delicate elsewhere in this program.

\subsection{The coefficients}

\begin{theorem}[All-rank local coefficients]
\label{thm:wickcoeffs}
The terminating conjugation gives exact shell projectors and a unique
all-order simple-branch formal recursion in $\Q[N, N^{-1}]$, with state
degree at most $s + 4r$. The two formal local gaps have exact all-rank
coefficients through $\beta^{-2}$, including $c_3^{-}$ and
$c_4^{\pm}$. The five corrections $c_0, \ldots, c_4$ are strictly
negative at every allowed rank, by positive shifted numerator
polynomials.
\end{theorem}
\gid{RESULT:LOCAL\_CLASS\_WICK\_RECURRENCE, RESULT:LOCAL\_CLASS\_WICK\_COEFFICIENTS, RESULT:LOCAL\_CLASS\_WICK\_SIGNS}

The construction applies to the vacuum, the first charge-even shell for
$N \ge 2$, and the first charge-odd shell for $N \ge 3$. It recovers the
previously archived $c_0$, $c_1$, $c_2$ formulas in both sectors and the
archived all-rank $c_3^{+}$, and additionally derives $c_3^{-}$ and both
all-rank $c_4$ formulas. The certificate carries full polynomial
eigen-equation residuals rather than only agreement with the archived
constants --- a distinction worth stating, because agreement with a
previous value is a weaker check than satisfying the equation.

The recovery is by an independent route: the archive engine uses
full-GUE cut-and-join and then subtracts the trace, while this one
applies $D$ directly in traceless coordinates.

\subsection{A companion algebraic identity}

\begin{theorem}[Shifted-power carrier symbols]
\label{thm:rsmr}
For the cubic Laurent operators with $S = L_{\downarrow} - 4I$,
\begin{equation}
  \sigma(R S^m R) = (q-4)^m\big(q e_2 + 3 e_3\big) - 4\,\Pi_m(q)\, e_2^2
\end{equation}
for every natural $m$, with $\Pi$ defined polynomially also at $q = 0$.
\end{theorem}
\gid{RESULT:RSM\_R\_CARRIER\_SYMBOL\_MASTER\_THEOREM}

\noindent This closes the word table of \S\ref{sec:sixth} at all powers,
and is formalized in Lean as an all-power operator statement.

\begin{scopenote}
Theorem~\ref{thm:wickcoeffs} concerns the \emph{formal} local class
oscillator only. No error bound for the compact-group eigenvalues is
claimed, and local spectral coefficients do not by themselves identify
the interacting Wilson spectrum, an interacting-volume estimate, or a
continuum source measure. Theorem~\ref{thm:rsmr} is an algebra identity
and carries no sixth-order dynamical population claim.
\end{scopenote}
